\chapter{Recitation 1: Complex Numbers and Induction}


\section*{How to use these notes}
These notes are intended for students in all fields of engineering. They combine definitions, guiding ideas, examples, and detailed solutions.
\tIdea{\textbf{Goals of this recitation.} We will learn to convert between Cartesian and polar forms, calculate powers and roots, and interpret operations in the complex plane. We will finish by writing a proof by induction, explicitly stating the range for which the claim holds.}
\tIdea{\textbf{A working habit.} In every solution, state what is given, what is required, which definition or claim you are using, and why each operation is valid. Correct calculations and clear explanations complement one another.}

\section{Complex numbers: Cartesian form}
\Deff{
The imaginary unit $i$ satisfies $i^{2}=-1$.

The expression $z=x+iy$ is a complex number, where $x,y\in\mathbb{R}$.
\\
\begin{enumerate}
    \item
If $ \operatorname{Re}(z)=x=0$, then $z$ is purely imaginary.
    \item
If $ \operatorname{Im}(z)=y=0$, then $z$ is real.

\end{enumerate}
}
\tExSol{Over $\mathbb C$, solve the equation $z^2-4z+13=0$. Write a factorization into linear factors over $\mathbb C$, and determine whether such a factorization exists over $\mathbb R$.
}{
The discriminant is $\Delta=(-4)^2-4\cdot1\cdot13=16-52=-36$. Therefore,
\[
z_{1,2}=\frac{4\pm\sqrt{-36}}2=\frac{4\pm6i}2=2\pm3i.
\]
It follows that
\[
z^2-4z+13=(z-(2+3i))(z-(2-3i)).
\]
This is a factorization over $\mathbb C$. Over $\mathbb R$, the discriminant is negative, so there are no real roots and no factorization into two real linear factors.

}

\Deff{
\begin{itemize}
\item \textbf{Conjugate:} For $z=x+iy$, we write $z^{*}:=x-iy$.
\item \textbf{Modulus:} $|z|=\sqrt{x^{2}+y^{2}}\in\mathbb{R}$.
\item \textbf{Multiplicative inverse:} For $z\neq0$: $z^{-1}=\dfrac{z^{*}}{|z|^{2}}$.
\end{itemize}
}
\tExSol{Prove: \textbf{a}. $ \operatorname{Re}z=\dfrac{z+z^{*}}{2},\ \operatorname{Im}z=\dfrac{z-z^{*}}{2i}$.
\textbf{b}. $|z|\ge \operatorname{Im}z$, $|z|\ge \operatorname{Re}z$ for every $z\in\mathbb{C}$.
\textbf{c}. \HebParen{Homework} $(z^{*})^{*}=z$.
\textbf{d}. \HebParen{Homework} If $z$ is real, then $z^{*}=z$.}{
Write $z=x+iy$. Then
\[
\frac{z+z^{*}}{2}=x=\operatorname{Re}z,\qquad
\frac{z-z^{*}}{2i}=y=\operatorname{Im}z,
\]
and
\[
|z|=\sqrt{x^{2}+y^{2}}\ge |x|\ge x=\operatorname{Re}z,\qquad
|z|\ge |y|\ge y=\operatorname{Im}z.
\]
\textbf{\HebParen{c}} $(z^{*})^{*}=(x-iy)^{*}=x+iy=z$.\\
\textbf{\HebParen{d}} If $z$ is real, then $y=0$, so $z^{*}=x-0i=x=z$.
}
\subsection{Useful properties}
\begin{enumerate}
\item $|z|=|z^{*}|$
\item $z\cdot z^{*}=|z|^{2}$ \quad\HebParen{Multiplication by the conjugate gives a real number}
\item $|zw|=|z|\,|w|$
\item $(zw)^{*}=z^{*}w^{*}$
\end{enumerate}

\tExSol{Prove that, for $z\neq0$: $(z^{-1})^{*}=(z^{*})^{-1}$.}{
Since $z\ne0$, we also have $z^{*}\ne0$ and $|z|^{2}>0$. By the formula for the multiplicative inverse,
\[
(z^{-1})^{*}=\left(\frac{z^{*}}{|z|^{2}}\right)^{*}
=\frac{z}{|z|^{2}},\qquad
(z^{*})^{-1}=\frac{(z^{*})^{*}}{|z^{*}|^{2}}
=\frac{z}{|z|^{2}}.
\]
Both sides equal the same expression, so $(z^{-1})^{*}=(z^{*})^{-1}$.
}

\section{Polar form}

\subsection{Converting between Cartesian and polar forms}
\Deff{

For every $z=x+iy\ne0$, there exist $r=|z|>0$ and an angle $\theta$ such that
\[
z=r(\cos\theta+i\sin\theta)=re^{i\theta},\qquad
x=r\cos\theta,\quad y=r\sin\theta.
\]
Conversely, $r=\sqrt{x^{2}+y^{2}}$. When $x\ne0$, we have $\tan\theta=y/x$, but this ratio alone does not determine the quadrant.
\textbf{Remarks:}
\begin{itemize}
\item The angle is determined up to the addition of $2\pi k$, where $k\in\mathbb{Z}$. For $z=0$, we have $r=0$, and the argument is undefined.
\item All angles in exponents $e^{i\theta}$ will be expressed in radians. Degrees will be used only for numerical illustration.
\end{itemize}
}

\tExSol{Convert $z=3+3i$ and $w=-4+4i$ to polar form. Then write $a=6e^{i\pi/3}$ in Cartesian form. In each polar form, choose an angle in the interval $[0,2\pi)$.
}{
For $z$, we have $|z|=\sqrt{3^2+3^2}=3\sqrt2$. The components are equal and positive, so the angle is $\pi/4$:
\[
z=3\sqrt2e^{i\pi/4}.
\]
For $w$, we have $|w|=\sqrt{(-4)^2+4^2}=4\sqrt2$. The point lies in the second quadrant, and the ratio of the magnitudes of its components is $1$, so
\[
w=4\sqrt2e^{3i\pi/4}.
\]
Finally, we use Euler's identity and the trigonometric values at the angle $\pi/3$:
\[
a=6\left(\cos\frac\pi3+i\sin\frac\pi3\right)
=6\left(\frac12+i\frac{\sqrt3}2\right)=3+3\sqrt3i.
\]
The angles are in radians. Both given numbers have positive modulus, so the issue of an argument at zero does not arise.

}

\Needspace{0.47\textheight}
\subsection{Arithmetic operations}
\CourseFigure{complex_plane}{Cartesian and polar forms in the same plane: the modulus is the distance from the origin, and the argument is the direction angle.}
\Deff{
\textbf{Multiplication:}\quad $(r_{1}e^{i\theta_{1}})(r_{2}e^{i\theta_{2}})=r_{1}r_{2}\,e^{i(\theta_{1}+\theta_{2})}$.

\textbf{Division} \HebParen{when the denominator is $\neq0$}:\quad
$\dfrac{r_{1}e^{i\theta_{1}}}{r_{2}e^{i\theta_{2}}}=\dfrac{r_{1}}{r_{2}}\,e^{i(\theta_{1}-\theta_{2})}$.

\textbf{Conjugate}:\quad $(re^{i\theta})^{*}=r\,e^{-i\theta}$.

\textbf{Periodicity:}\quad $re^{i(\theta+2\pi k)}=re^{i\theta}$,\ \ $re^{i(\theta+\pi)}=-re^{i\theta}$.
}
\tExSol{Let $z=1+i$ and $w=-4+4i$. Calculate $zw$ and $w/z$ in two ways: using polar form and using a Cartesian calculation. Check that the resulting moduli agree with the multiplication and division formulas.
}{
Write $z=\sqrt2e^{i\pi/4}$ and $w=4\sqrt2e^{3i\pi/4}$. Therefore,
\[
zw=8e^{i\pi}=-8,\qquad \frac wz=4e^{i\pi/2}=4i.
\]
A Cartesian calculation gives
\[
(1+i)(-4+4i)=-4+4i-4i+4i^2=-8,
\]
\[
\frac{-4+4i}{1+i}=\frac{(-4+4i)(1-i)}{(1+i)(1-i)}
=\frac{-4+8i+4}{2}=4i.
\]
Checking the moduli: $|z||w|=\sqrt2\cdot4\sqrt2=8=|-8|$,
and $|w|/|z|=4=|4i|$. The two calculation methods agree.

}

\subsection{Raising to a power}

\tExSol{Calculate $(1+i)^{12}$ using polar form. Verify the result using $(1+i)^4$.
}{
Since $1+i=\sqrt2e^{i\pi/4}$,
\[
(1+i)^{12}=(\sqrt2)^{12}e^{12i\pi/4}=2^6e^{3i\pi}=-64.
\]
To verify this, $(1+i)^2=2i$, so $(1+i)^4=(2i)^2=-4$.
Hence $(1+i)^{12}=((1+i)^4)^3=(-4)^3=-64$.

}

\subsection{Taking a complex root}
\Deff{
Let $0\neq z\in\mathbb{C}$, and let $n\ge1$ be an integer. The equation $w^{n}=z=re^{i\theta}$ has $n$ distinct roots:
\[
w_{j}=\sqrt[n]{r}\; e^{\,i\frac{\theta+2\pi j}{n}},\quad j=0,1,\ldots,n-1.
\]
\textbf{Geometric meaning:} The roots lie on a circle of radius $\sqrt[n]{r}$, with equal angular spacing $2\pi/n$. When $n\ge3$, they are the vertices of a regular polygon. For $n=1$ there is one root; for $n=2$ there are two opposite roots.
}
\tExSol{Find all solutions of $z^{4}=-i$.}{
Write $-i=e^{-i\pi/2}$. Therefore,
\[
z_{j}=e^{\,i\frac{-\pi/2+2\pi j}{4}}
= e^{\,i\frac{\pi(4j-1)}{8}},\quad j=0,1,2,3,
\]
that is, the angles are $-\pi/8,\ 3\pi/8,\ 7\pi/8,\ 11\pi/8$ \HebParen{at radius $1$}.
}

\tExSol{Let $n\ge2$. Find all $n$th roots of unity and prove that their sum is zero. What happens when $n=1$?}{
The roots are $w_{j}=e^{\,i\frac{2\pi j}{n}}$, $j=0,\dots,n-1$.
Write $w=e^{\,i\frac{2\pi}{n}}$. Then $w_{j}=w^{j}$, so
\[
\sum_{j=0}^{n-1} w^{j}=\frac{w^{n}-1}{w-1}=0,
\]
because $w^{n}=1$ and $w\ne1$ when $n\ge2$. For $n=1$, the only root is $1$, so the sum is $1$, not zero.
}

\tExSol{\HebParen{\textbf{Homework}} Consider $a_{n}x^{n}+a_{n-1}x^{n-1}+\cdots+a_{1}x+a_{0}=0$ with real coefficients.
Prove that if $z$ is a solution, then $z^{*}$ is also a solution.}{
We use the identities
$(z_{1}+\cdots+z_{m})^{*}=z_{1}^{*}+\cdots+z_{m}^{*}$,
$(z_{1}\cdots z_{m})^{*}=z_{1}^{*}\cdots z_{m}^{*}$,
and the fact that a real number $a$ satisfies $a^{*}=a$.
If $a_{0}+a_{1}z+\cdots+a_{n}z^{n}=0$, applying $^{*}$ gives
$a_{0}+a_{1}z^{*}+\cdots+a_{n}(z^{*})^{n}=0$.
Thus $z^{*}$ is a solution.
}

\Needspace{0.47\textheight}
\section{Further properties of complex numbers}
\CourseFigure{roots_unity}{Roots of unity are equally spaced in angle on the unit circle; for an order greater than 1, their sum is zero.}
\subsection{Multiplication by $i$}
Multiplication by $i$ rotates a complex number counterclockwise through an angle of \( \frac{\pi}{2} \).

\tExSol{
Prove that multiplication by \( i\) produces a rotation through an angle of \( \frac{\pi}{2} \).}
{

\textbf{Cartesian representation:}\\
Suppose \( z = a + bi \). Multiplication by $i$ gives
\[
z \cdot i = (a + bi)i = -b + ai.
\]
The vector $(a,b)$ has been replaced by the vector $(-b,a)$. This is exactly a counterclockwise rotation through an angle of $\pi/2$ about the origin: the new $x$ coordinate is $-b$, and the new $y$ coordinate is $a$.

\textbf{Polar representation:}\\
The case $z=0$ is immediate: the origin remains fixed. For a nonzero complex number, suppose \( z = r e^{i\theta} \), where \( r = |z| \) and \( \theta \) is the angle that \( z \) makes with the positive \( x \) axis. Multiplication by $i$ can be written as
\[
z \cdot i = r e^{i\theta} \cdot i = r e^{i\theta} \cdot e^{i\frac{\pi}{2}} = r e^{i(\theta + \frac{\pi}{2})}.
\]
Multiplication by $i$ increases \( \theta \) by the angle \( \frac{\pi}{2} \), meaning that the complex number \( z \) rotates through an angle of \( \frac{\pi}{2} \) about the origin, as required.
}
\subsection{Division in polar form}
\Deff{
If \( z_1 = r_1 e^{i\theta_1} \) and \( z_2 = r_2 e^{i\theta_2} \), and also \(r_2 \neq 0\), then
\[
\frac{z_1}{z_2} = \frac{r_1}{r_2} e^{i(\theta_1 - \theta_2)}.
\]}
\tExSol{
Calculate \( \frac{z_1}{z_2} \) for \( z_1 = 4e^{i\frac{\pi}{3}} \) and \( z_2 = 2e^{i\frac{\pi}{6}} \).}
{We calculate:
\[
\frac{z_1}{z_2} = \frac{4}{2} e^{i\left(\frac{\pi}{3} - \frac{\pi}{6}\right)} = 2 e^{i\frac{\pi}{6}}.
\]}
\tExSol{Let $z_1,z_2\ne0$. Prove that an argument of the quotient $z_1/z_2$ equals the difference of the arguments, up to a multiple of $2\pi$.}
{
Suppose \( z_1 = r_1 e^{i\theta_1} \) and \( z_2 = r_2 e^{i\theta_2} \), where $r_1,r_2>0$. Division gives
\[
\frac{z_1}{z_2} = \frac{r_1}{r_2} e^{i(\theta_1 - \theta_2)}.
\]
Hence $\theta_1-\theta_2$ is an argument of the quotient. If an argument in a fixed range is required, add a suitable multiple of $2\pi$.
}

\subsection{Choosing an argument in a fixed range}
In these notes, $\operatorname{Arg}(z)$ denotes the unique argument in the range $[0,2\pi)$, for $z\ne0$.
\tRem{ We cannot always write \( \text{Arg}(z_1 \cdot z_2) = \text{Arg}(z_1) + \text{Arg}(z_2) \) or \( \text{Arg}\left(\frac{z_1}{z_2}\right) = \text{Arg}(z_1) - \text{Arg}(z_2) \), because the result may fall outside the range \( [0, 2\pi) \).
}
\tExSol{
Calculate \( \text{Arg}(z_1 \cdot z_2) \) for \( z_1 = e^{i\frac{3\pi}{4}}, z_2 = e^{i\frac{5\pi}{4}} \).
}
{
First, calculate the product itself:
\[
z_1z_2=e^{i3\pi/4}e^{i5\pi/4}=e^{i2\pi}=1.
\]
The angle $2\pi$ is an argument of $1$, but it is not in the range $[0,2\pi)$. Therefore,
\[
\operatorname{Arg}(z_1z_2)=\operatorname{Arg}(1)=0,
\qquad
\operatorname{Arg}(z_1)+\operatorname{Arg}(z_2)=2\pi.
\]
The two numbers are not equal; they are congruent modulo $2\pi$.
}

\subsection{Calculating the argument and analyzing quadrants}
\Deff{
For a complex number \( z = a + bi\ne0 \), where $a\ne0$, every argument $\theta$ satisfies
\[
\tan(\theta) = \frac{b}{a}.
\]To determine \( \theta \) precisely, we must take into account the quadrant in which the complex number lies:
\begin{itemize}
    \item If \( a > 0, b > 0 \): the angle is \( \theta = \arctan\left(\frac{b}{a}\right) \).
    \item If \( a < 0, b > 0 \): add \( \pi \) to account for the second quadrant: \( \theta = \pi + \arctan\left(\frac{b}{a}\right) \).
    \item If \( a < 0, b < 0 \): add \( \pi \) for the third quadrant: \( \theta = \pi + \arctan\left(\frac{b}{a}\right) \).
    \item If \( a > 0, b < 0 \): add \( 2\pi \) to obtain a positive angle in the fourth quadrant: \( \theta = 2\pi + \arctan\left(\frac{b}{a}\right) \).
\end{itemize}
If $b=0$ and $a>0$, then $\operatorname{Arg}(z)=0$; if $b=0$ and $a<0$, then $\operatorname{Arg}(z)=\pi$.
If \( a = 0 \), the angle depends only on the value of \( b \):
\begin{itemize}
    \item If \( b > 0 \): \( \theta = \frac{\pi}{2} \).
    \item If \( b < 0 \): \( \theta = \frac{3\pi}{2} \).
\end{itemize}
}

\tExSol{Find the angle \( \theta \) for \( z = -1 + \sqrt{3}i \).
}
{
We are given \( a = -1, b = \sqrt{3} \). First, calculate
\[
\tan(\theta) = \frac{\sqrt{3}}{-1} = -\sqrt{3}.
\]
Since \( a < 0 \) and \( b > 0 \), the number lies in the second quadrant. Therefore,
\[
\theta = \pi - \frac{\pi}{3} = \frac{2\pi}{3}.
\]}
\tExSol{
Find the angle \( \theta \) for \( z = -5 - 5i \).
}
{
We are given \( a = -5, b = -5 \). First, calculate
\[
\tan(\theta) = \frac{b}{a} = \frac{-5}{-5} = 1.
\]
Since \( a < 0, b < 0 \), the number lies in the third quadrant. Therefore,
\[
\theta = \pi + \arctan(1) = \pi + \frac{\pi}{4} = \frac{5\pi}{4}.
\]
}



\section{The principle of induction}
\Deff{
\textbf{The principle of \HebParen{weak} induction:}\\
Let $n_0$ be a nonnegative integer, and let $P(n)$ be a statement defined for every integer $n\ge n_0$. Suppose that
\begin{enumerate}
    \item \textbf{Base case:} The statement $P(n_0)$ is true.
    \item \textbf{Induction step:} For every integer $n\ge n_0$, the assumption $P(n)$ implies the statement $P(n+1)$.
\end{enumerate}
Then $P(n)$ is true for every integer $n\ge n_0$.
}

\tExSol{
Prove by induction that, for every \(n\ge1\),
\[
1+2+\cdots+n=\frac{n(n+1)}{2}.
\]
}
{
\textbf{Defining the statement:} \(P(n):\; \sum_{i=1}^n i=\frac{n(n+1)}{2}\).\\
\textbf{Base case:} \(n=1\): the left-hand side \(=1\), and the right-hand side \(=\frac{1\cdot 2}{2}=1\). The statement is true.\\
\textbf{Induction hypothesis:} Suppose that, for some \(k\ge1\), we have \(\sum_{i=1}^{k} i=\frac{k(k+1)}{2}\).\\
\textbf{Induction step:}
\[
\sum_{i=1}^{k+1} i=\Big(\sum_{i=1}^{k} i\Big)+(k+1)=\frac{k(k+1)}{2}+(k+1)=\frac{(k+1)(k+2)}{2}.
\]
We have obtained the formula for \(k+1\).\\
\textbf{Conclusion:} The statement is true for every \(n\ge1\).
}
\tExSol{
Prove that, for every \(n\ge1\), we have \(7\,\mid\,\left(8^{n}-1\right)\).
}
{
\textbf{Defining the statement:} \(P(n):\; 7\mid (8^n-1)\).\\
\textbf{Base case:} \(n=1\): \(8^1-1=7\), and indeed \(7\mid 7\).\\
\textbf{Induction hypothesis:} Suppose that, for some \(k\ge1\), we have \(7\mid (8^k-1)\); that is, there exists \(m\in\mathbb{Z}\) such that \(8^k-1=7m\).\\
\textbf{Induction step:}
\[
8^{k+1}-1=8\cdot 8^k-1=8(8^k-1)+7.
\]
By the hypothesis, \(8^k-1=7m\), so
\[
8^{k+1}-1=8\cdot 7m+7=7(8m+1),
\]
and therefore \(7\mid (8^{k+1}-1)\).\\
\textbf{Conclusion:} For every \(n\ge1\), we have \(7\,\mid\,\left(8^{n}-1\right)\). \\
\textbf{A shorter argument:} \(8\equiv 1\pmod{7}\Rightarrow 8^n\equiv 1 \pmod{7}\).
}

\tExSol{
Prove that, for every \(n\ge 4\),
\[
2^{\,n}\ge n^{2}.
\]
}
{
We prove this by induction on \(n\).\\
\textbf{Base case:} \(n=4\): \(2^4=16=4^2\). The statement is true.\\
\textbf{Induction hypothesis:} Suppose \(2^k\ge k^2\) for \(k\ge 4\).\\
\textbf{Induction step:}
\[
2^{k+1}=2\cdot 2^k \;\ge\; 2k^2.
\]
It suffices to show that \(2k^2\ge (k+1)^2 \iff k^2-2k-1\ge 0\).
For \(k\ge 4\), we have \(k^2-2k-1>0\), so \(2^{k+1}\ge (k+1)^2\).\\
Hence, by induction, \(2^{n}\ge n^{2}\) for every \(n\ge 4\).
}

\section{Supplementary exercises: geometry and signal calibration}
\tExSol{Find and describe geometrically all $z\in\mathbb{C}$ satisfying $|z-1|=|z+1|$. Which of them also satisfy $|z|=2$?}{
Write $z=x+iy$. Both sides are nonnegative, so squaring is valid:
\[
|z-1|^{2}=(x-1)^{2}+y^{2},\qquad
|z+1|^{2}=(x+1)^{2}+y^{2}.
\]
It follows that
\[
(x-1)^2+y^2=(x+1)^2+y^2
\iff x^2-2x+1=x^2+2x+1
\iff x=0.
\]
This is the imaginary axis: the set of points equidistant from the two points $1,-1$. The additional condition gives $y^2=4$, so the intersection consists of $z=2i$ or $z=-2i$.
}
\tExSol{In a simple model of a signal with two components, a calibration system performs the operation $T(z)=(1+i)z$. What is the effect on the magnitude and angle? Calculate $T(1-i)$, and find an input for which the output is $2i$.}{
Since $1+i=\sqrt2e^{i\pi/4}$, for every $z=re^{i\theta}\ne0$, we have $T(z)=\sqrt2\,r\,e^{i(\theta+\pi/4)}$.
Thus, the magnitude is multiplied by $\sqrt2$, and the direction rotates counterclockwise through an angle of $\pi/4$. We also have $T(0)=0$.
A direct calculation gives $T(1-i)=(1+i)(1-i)=1-i^2=2$.
To obtain $2i$, solve $(1+i)z=2i$:
\[
z=\frac{2i}{1+i}=\frac{2i(1-i)}{(1+i)(1-i)}
=\frac{2i+2}{2}=1+i.
\]
Check: $(1+i)^2=1+2i+i^2=2i$.
}
